\documentclass[10pt,a4paper]{amsart}
\usepackage{amsmath,amssymb,amsthm,graphicx}
\usepackage{verbatim}
\newenvironment{reference}{\comment}{\endcomment}
%\newenvironment{reference}{}{}
\newenvironment{slogan}{\comment}{\endcomment}
\newenvironment{history}{\comment}{\endcomment}

% For commutative diagrams we use Xy-pic
\usepackage[all]{xy}

% We use 2cell for 2-commutative diagrams.
\xyoption{2cell}
% Theorem environments.
%
\theoremstyle{plain}
\newtheorem{theorem}[subsection]{Theorem}
\newtheorem{proposition}[subsection]{Proposition}
\newtheorem{lemma}[subsection]{Lemma}

\theoremstyle{definition}
\newtheorem{definition}[subsection]{Definition}
\newtheorem{example}[subsection]{Example}
\newtheorem{exercise}[subsection]{Exercise}
\newtheorem{situation}[subsection]{Situation}

\theoremstyle{remark}
\newtheorem{remark}[subsection]{Remark}
\newtheorem{remarks}[subsection]{Remarks}

\numberwithin{equation}{subsection}

% Macros
%
\def\lim{\mathop{\mathrm{lim}}\nolimits}
\def\colim{\mathop{\mathrm{colim}}\nolimits}
\def\Spec{\mathop{\mathrm{Spec}}}
\def\Hom{\mathop{\mathrm{Hom}}\nolimits}
\def\Ext{\mathop{\mathrm{Ext}}\nolimits}
\def\SheafHom{\mathop{\mathcal{H}\!\mathit{om}}\nolimits}
\def\SheafExt{\mathop{\mathcal{E}\!\mathit{xt}}\nolimits}
\def\Sch{\mathit{Sch}}
\def\Mor{\mathop{\mathrm{Mor}}\nolimits}
\def\Ob{\mathop{\mathrm{Ob}}\nolimits}
\def\Sh{\mathop{\mathit{Sh}}\nolimits}
\def\NL{\mathop{N\!L}\nolimits}
\def\CH{\mathop{\mathrm{CH}}\nolimits}
\def\proetale{{pro\text{-}\acute{e}tale}}
\def\etale{{\acute{e}tale}}
\def\QCoh{\mathit{QCoh}}
\def\Ker{\mathop{\mathrm{Ker}}}
\def\Im{\mathop{\mathrm{Im}}}
\def\Coker{\mathop{\mathrm{Coker}}}
\def\Coim{\mathop{\mathrm{Coim}}}

% Boxtimes
%
\DeclareMathSymbol{\boxtimes}{\mathbin}{AMSa}{"02}

%
% Macros for moduli stacks/spaces
%
\def\QCohstack{\mathcal{QC}\!\mathit{oh}}
\def\Cohstack{\mathcal{C}\!\mathit{oh}}
\def\Spacesstack{\mathcal{S}\!\mathit{paces}}
\def\Quotfunctor{\mathrm{Quot}}
\def\Hilbfunctor{\mathrm{Hilb}}
\def\Curvesstack{\mathcal{C}\!\mathit{urves}}
\def\Polarizedstack{\mathcal{P}\!\mathit{olarized}}
\def\Complexesstack{\mathcal{C}\!\mathit{omplexes}}
% \Pic is the operator that assigns to X its picard group, usage \Pic(X)
% \Picardstack_{X/B} denotes the Picard stack of X over B
% \Picardfunctor_{X/B} denotes the Picard functor of X over B
\def\Pic{\mathop{\mathrm{Pic}}\nolimits}
\def\Picardstack{\mathcal{P}\!\mathit{ic}}
\def\Picardfunctor{\mathrm{Pic}}
\def\Deformationcategory{\mathcal{D}\!\mathit{ef}}
\usepackage{cleveref}
\crefname{originallemma}{original lemma}{original lemmas}
\Crefname{originallemma}{Original lemma}{Original lemmas}
\makeatletter
\newlabel{duality-lemma-sheaf-with-exact-support-ext}{{0A77}{source}{}{sourcepointer-1}{}}
\newlabel{duality-lemma-sheaf-with-exact-support-ext@cref}{{[originallemma][0][]0A77}{[1][0][]source}{}{}{}}
\makeatother
\crefname{originallemma}{original lemma}{original lemmas}
\Crefname{originallemma}{Original lemma}{Original lemmas}
\makeatletter
\newlabel{duality-lemma-twisted-inverse-image}{{0A9E}{source}{}{sourcepointer-2}{}}
\newlabel{duality-lemma-twisted-inverse-image@cref}{{[originallemma][0][]0A9E}{[1][0][]source}{}{}{}}
\makeatother
\crefname{originallemma}{original lemma}{original lemmas}
\Crefname{originallemma}{Original lemma}{Original lemmas}
\makeatletter
\newlabel{duality-lemma-twisted-inverse-image-closed}{{0A9X}{source}{}{sourcepointer-3}{}}
\newlabel{duality-lemma-twisted-inverse-image-closed@cref}{{[originallemma][0][]0A9X}{[1][0][]source}{}{}{}}
\makeatother
\crefname{originalreference}{original reference}{original references}
\Crefname{originalreference}{Original reference}{Original references}
\makeatletter
\newlabel{lemma-duality-dim-1}{{0BS2}{source}{}{sourcepointer-4}{}}
\newlabel{lemma-duality-dim-1@cref}{{[originalreference][0][]0BS2}{[1][0][]source}{}{}{}}
\makeatother
\title{Dualizing complex and sheaf of a closed subscheme by RSheafHom against the ambient ones}
\author{The Stacks Project Authors}
\date{stacks-project:curves.tex:local-snapshot-20261012; GNU FDL 1.2 or later}
\begin{document}
\maketitle
\noindent\textit{Editorial note.} All delivered mathematical text is the author's own, copied from the pinned Stacks Project chapter file curves.tex; no statement or proof step was rewritten, reordered or removed.\par\medskip
\noindent\textit{Editorial note.} The counted proof is the author's complete argument: the adjoint composition, the RSheafHom identification, the H\textasciicircum{}\{-1\} extraction, and the pushforward upgrade are all retained verbatim.\par\medskip
\noindent\textit{Editorial note.} The label lemma-closed-immersion-dim-1-CM names the counted statement, and the proof environment immediately following it in the source is delivered verbatim as its proof.\par\medskip
\noindent\textit{Editorial note.} The counted statement is carried by the author's own lemma environment in the chapter source.\par\medskip
\noindent\textit{Allowed standard prerequisites (editorial).}\begin{itemize}
\item Twisted inverse images as right adjoints. Duality for Schemes, Lemma duality-lemma-twisted-inverse-image and duality-lemma-twisted-inverse-image-closed. Provides the adjoints a, b, c of f, g, i, the composition b = c composed with a from uniqueness of right adjoints, and the RSheafHom description of c.
\item The curve duality framework. Same chapter, Lemma lemma-duality-dim-1. Fixes omega\_X\textasciicircum{}bullet = a(O\_Spec(k)) and omega\_X = H\textasciicircum{}\{-1\}, the objects related by the lemma.
\item Sheaf Hom against ideals. Duality for Schemes, Lemma duality-lemma-sheaf-with-exact-support-ext. Upgrades the H\textasciicircum{}\{-1\} identification to i\_*omega\_Y = SheafHom\_\{O\_X\}(i\_*O\_Y, omega\_X).
\end{itemize}\medskip
\ifdefined\hypertarget\hypertarget{sourcepointer-1}{}\fi
\noindent\textit{Source reference.} Original lemma 0A77: sources/stacks-chow/tags.txt, source lines 12870--12870. Prerequisite pointer invoked at proof line(s) 1008; it names a standard fact of a sibling chapter that the counted argument uses, and every step of the counted proof built on it is retained verbatim in the delivered body.\par\medskip
\ifdefined\hypertarget\hypertarget{sourcepointer-2}{}\fi
\noindent\textit{Source reference.} Original lemma 0A9E: sources/stacks-chow/tags.txt, source lines 12947--12947. Prerequisite pointer invoked at proof line(s) 992; it names a standard fact of a sibling chapter that the counted argument uses, and every step of the counted proof built on it is retained verbatim in the delivered body.\par\medskip
\ifdefined\hypertarget\hypertarget{sourcepointer-3}{}\fi
\noindent\textit{Source reference.} Original lemma 0A9X: sources/stacks-chow/tags.txt, source lines 12965--12965. Prerequisite pointer invoked at proof line(s) 1001; it names a standard fact of a sibling chapter that the counted argument uses, and every step of the counted proof built on it is retained verbatim in the delivered body.\par\medskip
\ifdefined\hypertarget\hypertarget{sourcepointer-4}{}\fi
\noindent\textit{Source reference.} Original reference 0BS2: sources/stacks-curves/pinned-original-source.tex, source lines 796--796. Prerequisite pointer invoked at proof line(s) 980; it names a standard fact of the same chapter that the counted argument uses, and every step of the counted proof built on it is retained verbatim in the delivered body.\par\medskip
\begin{lemma}
\label{lemma-closed-immersion-dim-1-CM}
Let $X$ be a proper scheme of dimension $\leq 1$ over a field $k$.
Let $i : Y \to X$ be a closed immersion.
Let $\omega_X^\bullet$, $\omega_X$, $\omega_Y^\bullet$, $\omega_Y$
be as in Lemma \ref{lemma-duality-dim-1}. Then
\begin{enumerate}
\item $\omega_Y^\bullet = R\SheafHom(\mathcal{O}_Y, \omega_X^\bullet)$,
\item $\omega_Y = \SheafHom(\mathcal{O}_Y, \omega_X)$ and
$i_*\omega_Y = \SheafHom_{\mathcal{O}_X}(i_*\mathcal{O}_Y, \omega_X)$.
\end{enumerate}
\end{lemma}

\begin{proof}
Denote $g : Y \to \Spec(k)$ and $f : X \to \Spec(k)$ the structure morphisms.
Then $g = f \circ i$. Denote $a, b, c$ the right adjoint of
Duality for Schemes, Lemma
\ref{duality-lemma-twisted-inverse-image} for $f, g, i$.
Then $b = c \circ a$ by uniqueness of right adjoints
and because $Rg_* = Rf_* \circ Ri_*$.
In the proof of Lemma \ref{lemma-duality-dim-1}
we set
$\omega_X^\bullet = a(\mathcal{O}_{\Spec(k)})$ and
$\omega_Y^\bullet = b(\mathcal{O}_{\Spec(k)})$.
Hence $\omega_Y^\bullet = c(\omega_X^\bullet)$
which implies (1) by Duality for Schemes, Lemma
\ref{duality-lemma-twisted-inverse-image-closed}.
Since $\omega_X = H^{-1}(\omega_X^\bullet)$ and
$\omega_Y = H^{-1}(\omega_Y^\bullet)$ we conclude that
$\omega_Y = \SheafHom(\mathcal{O}_Y, \omega_X)$.
This implies
$i_*\omega_Y = \SheafHom_{\mathcal{O}_X}(i_*\mathcal{O}_Y, \omega_X)$
by Duality for Schemes, Lemma
\ref{duality-lemma-sheaf-with-exact-support-ext}.
\end{proof}

\end{document}
